\ficha{Salomão (2017), Controvérsias monetárias no Brasil imperial}{salomao2017}

\begin{identificacao}
SALOMÃO, Ivan Colangelo. Controvérsias monetárias no Brasil imperial e suas
influências na formação do pensamento desenvolvimentista brasileiro.
\emph{Desenvolvimento em Questão}, Ijuí, ano 15, n. 41, p. 6-23, out./dez.
2017. DOI 10.21527/2237-6453.2017.41.6-23.

Artigo de periódico, história do pensamento econômico brasileiro, 18 páginas,
em português. Lido integralmente. Autor: orientador da dissertação.
\end{identificacao}

\bloco{Problema e tese}
O artigo pergunta de onde vêm as ideias que sustentam a política econômica
ativa adotada no Brasil a partir de 1930, e responde que uma de suas raízes
está no papelismo do debate monetário imperial. A tese tem duas partes. A
primeira é de filiação: a controvérsia brasileira entre metalistas e papelistas
foi importada das escolas britânicas, mas foi reprocessada para refletir
particularidades da economia nacional, e é isso que o autor chama de
tropicalização. A segunda é de continuidade: o papelismo é um dos três
componentes do ``núcleo duro'' do desenvolvimentismo, ao lado do nacionalismo e
da industrialização, na definição de Fonseca (2004) que o artigo adota (p. 8).

\bloco{Argumento}
\begin{enumerate}[leftmargin=*,nosep]
\item O desenvolvimentismo como fato é do século 20, mas o fato não exige
concomitância com o corpo teórico: elementos dele já se pronunciavam no meio
jornalístico, militar e político do fim do século 19 (p. 8). Entre as três
correntes formadoras, o papelismo é a mais controversa e a de formulação mais
elaborada, porque precisava recorrer aos clássicos para se contrapor aos
metalistas (nota 3, p. 8). O eixo do debate imperial é a conversibilidade:
políticos de formação liberal abraçavam o padrão-ouro e a estabilidade cambial,
enquanto os ligados às atividades produtivas deslocavam o eixo para o nível de
liquidez e pediam atenção também à taxa de juros (p. 10).

\item Sobre o modo da importação, o texto opõe duas leituras e escolhe uma.
Furtado (1982, p. 160), citado em bloco, vê mimetismo, doutrina europeia aceita
sem confronto com a realidade. Saes (1983), Gremaud (1997) e Gambi (2011)
sustentam que a ideologia já chegava processada (p. 11). O resumo afirma a
tropicalização como resultado, mas o corpo do texto hesita e escreve
``tropicalizado ou não'' antes de passar à argumentação original (p. 11).

\item Reconstrói duas controvérsias europeias. A primeira opõe
\emph{bullionists} a \emph{anti-bullionists}: os primeiros, com Thornton e
Wheatley apoiados em Smith, Ricardo e Mill, atribuem valor intrínseco à moeda e
fazem da conversibilidade instrumento contra a inflação; os segundos partem das
``necessidades do mercado'', tratam a moeda como signo e atribuem a inflação a
causas alheias à política monetária expansiva (p. 13). A segunda opõe
\emph{currency} e \emph{banking school}, com um ponto de convergência: ambas
tomam a conversibilidade-ouro como regra necessária, o que Fonseca e Mollo
(2012, p. 9) chamam de ortodoxização da discussão (p. 14). A divergência é
sobre controles quantitativos de curto prazo, com monopólio emissor e expansão
só por entrada líquida de ouro de um lado, e autorregulação, \emph{real bills
doctrine} e pluralidade emissora de outro, esta com Bosanquet, Torrens, Boase e
sobretudo Attwood, da escola de Birmingham (p. 14).

\item A transposição brasileira mapeia os dois campos. Os metalistas defendem
padrão metálico e monopólio emissor (p. 15), e o argumento decisivo é o
ajustamento automático inerente à moeda metálica, ausente na fiduciária (p.
15). Os papelistas defendem inconversibilidade e pluralidade emissora, e querem
a emissão a cargo dos bancos e não do Tesouro, por inadequação temporal e
geográfica do monopólio público às demandas do mercado (p. 16 e p. 17).
Sustentam que a taxa de câmbio é determinada pelo balanço de pagamentos, e não
pela expansão do papel, defesa contra a acusação feita a Rui Barbosa a partir
de 1890 (p. 16).

\item O ponto fraco reconhecido é o controle da emissão. Abandonada a
conversibilidade, some o autocontrole do sistema bancário, e resta a prudência
dos gestores públicos ou a concorrência entre bancos, principal argumento
papelista pela pluralidade emissora (p. 18). O próprio Mauá formula a
dificuldade como a questão mais difícil de resolver (p. 18). Os episódios que
ancoram a corrente são a reforma de Souza Franco nos anos 1850 e seu discurso
de 1857 sobre o juro alto como obstáculo à produção (p. 17 e p. 18), e o
panfleto \emph{O Meio-Circulante no Brasil}, de Mauá, de 1878 (p. 18).

\item O fecho liga papelismo e desenvolvimentismo pela admissão do crédito e do
déficit público como instrumentos, e pela ideia de um Estado responsável pelo
crescimento e atuante como agente anticíclico, na formulação de Fonseca (2004,
p. 11) que o autor cita (p. 20). Rui Barbosa é o ensaio embrionário, quatro
décadas antes de Vargas, e o autor da expressão ``desenvolvimento nacional''
nos anos 1890 (p. 20). A confluência plena das três correntes, somado o
positivismo, só ocorre em meados do século 20 (p. 20 e p. 21).
\end{enumerate}

\chave{já era processada de modo a refletir as particularidades da economia
brasileira}{p. 11}

\chave{ambos assumiam a conversibilidade-ouro da moeda como regra necessária}{p. 14}

\chave{a experiência, e não uma teoria, é que deveria balizar o caminho mais
apropriado a ser seguido}{p. 16}

\bloco{Conceitos e definições operacionais}
Padrão-ouro: taxa de câmbio fixa e totalmente conversível em ouro ou na moeda
de transação internacional, a libra esterlina (p. 12). Metalismo: padrão
monetário metálico mais monopólio emissor, com o equilíbrio cambial automático
como virtude central (p. 15). Papelismo: inconversibilidade mais pluralidade
emissora, tendo como variável principal a oferta monetária condizente com o
ânimo dos negócios (p. 16). \emph{Real bills doctrine}, ou lei do refluxo: os
empréstimos, uma vez saldados, retornam ao sistema bancário, de modo que a
emissão não gera excesso de meio circulante (p. 14 e p. 17). Ortodoxização da
discussão: o efeito de ambas as escolas inglesas aceitarem a conversibilidade
como regra, o que estreita o desacordo aos controles quantitativos (p. 14).
Tropicalização: reprocessamento das ideias importadas segundo as
particularidades da economia local (p. 11).

\bloco{Evidência e método}
História das ideias por reconstrução doutrinária, sem série estatística nem
contrafactual. A base é literatura secundária: Fonseca, Fonseca e Mollo,
Gremaud, Saes, Saéz, Franco e Furtado. As falas dos protagonistas imperiais
entram quase todas por \emph{apud}: Souza Franco e seu discurso de 1857 via
Franco (1983, p. 83 e p. 91), Mauá via Fernandes (1974, p. 26) e via Barbosa
(1891, p. 258). Nenhum arquivo, anais parlamentares ou jornal é lido
diretamente. O período coberto vai do Segundo Reinado a 1890, com projeção
sobre 1930.

\bloco{Agentes e vertentes}
Metalistas brasileiros: Joaquim José Rodrigues Torres, Francisco Belisário
Soares de Sousa, Francisco de Sales Torres Homem e Joaquim Duarte Murtinho (p.
16). Papelistas: Bernardo de Sousa Franco, Irineu Evangelista de Sousa, o
Visconde de Mauá, e ainda Teixeira Júnior, Ouro Preto, João Alfredo, Lafayette
Rodrigues Pereira e sobretudo Rui Barbosa (nota 6, p. 16). Do lado europeu,
Thornton e Wheatley entre os bulionistas, apoiados em Smith, Ricardo e Mill (p.
13); Hume e sobretudo Ricardo comparecem antes como raiz doutrinária do polo
ortodoxo (p. 12); Bosanquet, Torrens, Boase e Attwood pela \emph{banking
school} (p. 14). Sobre bases sociais o autor situa os papelistas entre industriais,
produtores rurais e comerciantes (p. 16), e registra em nota a controvérsia
entre Faoro (2001) e Prado (2003) sobre a correspondência entre partidos e
interesses (nota 8, p. 19).

\bloco{Críticas e limites}
O recorte para em 1890 e não trata da Caixa de Conversão, de Taubaté nem do
período 1906-1914. A continuidade que o artigo afirma vai do papelismo imperial
a 1930 e salta a Primeira República, justamente onde o capítulo precisa
trabalhar. Há tensão interna sobre o estatuto teórico do papelismo: a nota 3
(p. 8) o chama de formulação rebuscada, e o corpo do texto diz que lhe faltava
corpo teórico de mesma envergadura, restando a razão prática (p. 16). A
filiação bulionista é imprecisa ao arrolar Mill entre os pais invocados por
Thornton e Wheatley, autores anteriores a ele (p. 13). A caracterização social
também vacila: produtores rurais aparecem entre os papelistas no corpo (p. 16)
e os representantes da lavoura agroexportadora entre os metalistas na nota 8
(p. 19). O contraponto de Franco (2008, p. 2), para quem as ideias se formam
como racionalização \emph{a posteriori} das medidas, é citado na introdução (p.
9) e não é enfrentado depois. Por fim, ler o século 19 a partir de 1930 impõe
risco de teleologia, que o autor mitiga ao negar que o desenvolvimentismo
surgisse já plenamente configurado (p. 20), sem eliminá-lo.

\begin{dialogo}
\item Sobrevivência do vocabulário escolar. Procurar nos quatro jornais menção
nominal a Ricardo, Peel, Attwood, \emph{banking school}, \emph{currency},
escola inglesa, lei do refluxo. Testa se em 1906 a autoridade doutrinária
inglesa ainda é invocada por nome, como no debate imperial reconstruído aqui,
ou se o argumento já circula sem pedigree declarado.

\item Monopólio contra pluralidade emissora. O artigo faz desse eixo o divisor
dos campos imperiais (p. 15 e p. 16). Procurar se em 1906 os artigos discutem
quem emite, Tesouro, Banco do Brasil ou a própria Caixa, e se ``liberdade
bancária'' ainda aparece. Testa a hipótese da dissertação de que o eixo
imperial não organiza mais o debate de 1906, que se dá entre valorização e
estabilidade a taxa nova.

\item Causalidade do câmbio e necessidades do comércio. O argumento papelista
é que o câmbio responde ao balanço de pagamentos, e não à emissão (p. 16), e
que a moeda nasce das ``necessidades do comércio'' (p. 17). Procurar quem
atribui a alta do mil-réis ao café e aos empréstimos, e não à Caixa, e procurar
equivalentes de época como escassez de numerário na defesa do limite de
emissão. Testa se a gramática da \emph{real bills doctrine} sustenta o campo
expansionista de 1906.

\item Rui Barbosa como pivô. O artigo o coloca entre os mais destacados
próceres papelistas e como ensaio desenvolvimentista embrionário (p. 10 e p.
20). Procurar
como os jornais de 1906-1914 o invocam, como autoridade ou como advertência do
Encilhamento. Testa se o papelismo já era vocabulário desacreditado quando a
Caixa foi criada.
\end{dialogo}

\usonocapitulo{Parte III, como texto de referência e como interlocutor a ser
posicionado. Serve a três afirmações. Primeira, que as categorias do debate
brasileiro vêm das duas controvérsias inglesas e chegam reprocessadas, tese da
tropicalização com literatura nomeada, Saes (1983), Gremaud (1997) e Gambi
(2011) contra o mimetismo de Furtado (1982). Segunda, que a ortodoxização
apontada por Fonseca e Mollo (2012) estreita o desacordo inglês aos controles
quantitativos, o que dá critério para medir o que atravessou o Atlântico.
Terceira, que os agentes imperiais dos dois campos podem ser nomeados, com
Murtinho entre os metalistas e Rui Barbosa entre os papelistas. O capítulo
precisa se demarcar em um ponto: o artigo trata o par metalismo e papelismo
como divisor vigente e projeta a linha direto sobre 1930, enquanto a
dissertação sustenta que em 1906 esse par já não divide os campos. Daí o uso,
caracterizar a herança de vocabulário, não classificar os agentes da Caixa.}
